\chapter{K近邻方法}
\chaptergoal{根据训练样本的距离直观解释分类与回归；正确理解邻居数与特征尺度；使用 sklearn 比较不同距离设置和邻居数；说明为何高维空间中的“相近”需要特别小心。}

K近邻（$k$-nearest neighbors, KNN）是一种非常直观的预测方法：面对一条新记录，先寻找训练数据中最相近的若干条，再根据邻居的已知目标作出预测。它不需要在训练阶段寻找线性系数或树的阈值，因此适合用来帮助学生理解“算法各有不同的学习机制”。
\section{相似样本与距离}

\textbf{邻居}不是对象在地理位置上相邻，而是指记录的特征值按所选距离计算相近。把两个输入分别作为横、纵坐标，一条记录就对应平面上的一个点；“寻找邻居”就是寻找离待预测点近的训练点。欧氏距离是平面中直线长度的推广，下面给出可以手算的例子。

对于两个数值特征，最常见的欧氏距离是 $d(\mathbf x,\mathbf z)=\sqrt{\sum_j(x_j-z_j)^2}$。若某个输入以千计，另一个只在0到1之间变化，前者的数值差异可能主导距离，即使两者在科学问题中的意义相近。因此，选择距离不仅是数学操作，还包含对尺度、测量单位和变量重要性的判断。

在物理实验中，一个特征是频率（Hz），另一个是流速（m/s），不经处理直接计算欧氏距离可能缺乏物理依据。标准化能够缓解单纯数值量级差异，但也不是唯一选择；如果已知某些变量的测量不确定性，距离或权重设计可能还要结合误差模型。KNN 的距离定义始终应与任务背景一起说明。

\subsection{距离不是对象固有的唯一属性}
K近邻法（K-nearest neighbors，KNN）依据所选特征中的邻近关系进行预测。首先必须规定怎样度量距离。二维欧氏距离为$\sqrt{(x_1-z_1)^2+(x_2-z_2)^2}$，推广到$p$维就是各列差平方的和再开平方；曼哈顿距离为各列差绝对值之和。二者都可以衡量分离程度，但形成的“附近区域”形状不同。

欧氏距离对应圆或球形邻域，曼哈顿距离在二维中对应菱形邻域。使用哪一种，应当服务于数据表示与任务，而不是因为名称熟悉。在高维测量中，字段单位、噪声和冗余都会改变距离含义。

\worked{先在纸上找邻居}
查询点$q=(0,0)$，三个训练点A=$(1,0)$、B=$(0,2)$、C=$(2,2)$，标签分别为0、1、1。欧氏距离依次为1、2、$\sqrt8\approx2.828$，所以最近一个邻居是A，三个邻居则包括全部记录。若第一维乘以10，距离变成10、2、$\sqrt{404}\approx20.1$，最近邻变成B。仅改变数值尺度就改变了预测依据。

这个构造不是建议把某一列任意放大，而是说明距离包含隐含权重。标准化、物理单位选择和测量误差处理，都是在决定哪些差异更重要。程序返回“最近”之前，研究者先决定了按什么标准近。

\subsection{加权距离和测量误差}
可以使用$d^2(x,z)=\sum_j (x_j-z_j)^2/s_j^2$，其中$s_j$是规定的比较尺度。用训练标准差时，这等价于先标准化再算欧氏距离；用已知测量精度时，则赋予精度较高的差异更高权重。若误差有相关性，单独按每列缩放仍不能完整表达，需要更一般的协方差结构。

本课程先使用训练标准化，但应保留对其物理含义的提问。两个波段的颜色差可能包含共同误差，两个能量特征可能重复描述同一信号；把它们当作独立方向会重复计权。第13章PCA和第16章相关特征分析将进一步讨论。

\section{K近邻分类与回归}

下面五条记录为人工构造，不对应真实观测。待预测点是 $\mathbf q=(2,2)$：
\begin{table}[htbp]\centering
\begin{tabular}{ccccc}\toprule
记录 & $x_1$ & $x_2$ & 已知类别 & 到$\mathbf q$的距离\\\midrule
A & 1 & 2 & 0 & $1$\\
B & 2 & 3 & 0 & $1$\\
C & 3 & 3 & 1 & $\sqrt2$\\
D & 5 & 2 & 1 & $3$\\
E & 5 & 5 & 1 & $\sqrt{18}$\\\bottomrule
\end{tabular}
\caption{用于邻居搜索的人工记录。记录名A、B等不是类别名。}
\end{table}
以C为例，距离是 $\sqrt{(3-2)^2+(3-2)^2}=\sqrt2$。将五个距离排序，最近的三个是A、B、C。其中类别0有两票，类别1有一票，所以 $k=3$、均匀投票时预测为类别0。我们没有拟合一条直线，而是使用了附近已有记录的答案。

\begin{lstlisting}[caption={核对三近邻的预测}]
import numpy as np
from sklearn.neighbors import KNeighborsClassifier
X = np.array([[1,2], [2,3], [3,3], [5,2], [5,5]])
y = np.array([0,0,1,1,1])
knn = KNeighborsClassifier(n_neighbors=3).fit(X, y)
print(knn.predict([[2,2]]))
print(knn.kneighbors([[2,2]])[0])
\end{lstlisting}
结果应为类别0及距离 $1,1,\sqrt2$。距离并列时，返回记录的先后可能依实现而异；如果并列发生在第$k$名边界，样本顺序也可能影响最终邻居集合。

分类时，模型先找出距离最小的 $k$ 个训练样本，根据其类别投票，或对较近的样本赋予更大权重。回归时，可取邻居目标值的均值或距离加权均值。例如有三名邻居目标值为10、12、15，均匀加权的预测是 $37/3$。这是一个人为设定的手算例子，不对应某种实际物理单位。

KNN 的 \texttt{fit()} 主要保存和组织训练数据，预测时才进行相似性搜索；它不应被描述为“梯度下降学习了参数”。邻居的选择在面对新样本时仍然可能非常耗时，尤其当训练集很大、特征维度较高时。

\subsection{训练阶段保存什么，预测阶段计算什么}
KNN的拟合主要保存或组织训练样本以及对应目标，可能建立便于检索的数据结构。预测时，为每个查询点寻找最近的$k$条训练记录，再按规定规则汇总目标。这里没有要通过梯度下降估计的一组线性权重。所谓“懒惰学习”强调主要工作发生在查询阶段，不表示它不使用训练数据。

分类的均匀权重规则是多数投票；类别概率可以用邻居中各类的比例作估计。回归的均匀规则是目标平均值。它们使用相同的邻居检索，却在汇总目标时不同。对类别代码0、1、2取普通平均会生成没有类别含义的1.4，不能把分类规则直接写成回归平均。

\worked{多数投票与距离加权投票可以不同}
沿用上节A、B、C的距离和标签，$k=1$预测0，$k=3$均匀投票预测1。若使用距离倒数作权重，类别0总权重为1，类别1总权重为$1/2+1/\sqrt8\approx0.8536$，于是预测0。两个结果都由清楚的规则得到；改变规则就改变了“附近多数”的含义。

若查询点恰好与训练点重合，距离倒数形式出现除零问题。实际实现必须规定处理方式，例如优先使用距离为零的记录，而不是直接产生无穷数。对真实应用，还应检查重合是否来自重复对象、复制记录或恰好相同的测量特征。

\subsection{回归平均为什么能被理解为局部常数拟合}
固定一个查询点及其$k$个邻居，在这些邻居的目标上寻找一个常数$c$，最小化$\sum_{i\in N_k}(y_i-c)^2$，解为邻居均值。因此KNN回归可以看作每个查询位置重新建立一个局部常数模型。使用不同距离权重时，最优常数变成相应加权平均。

例如三个邻居目标为10、12、20，均匀预测为14；若距离为1、2、4，权重为1、1/2、1/4，则预测为$(10+6+5)/(1+1/2+1/4)=12$。读者应同时报告邻居集合与权重，才完整说明预测来源。

\section{邻居数和预测规则}

\textbf{均匀投票}使每个入选邻居贡献一票；\textbf{距离加权}则让较近邻居贡献更多。例如三个邻居到查询点的距离为 $1,2,4$，用距离倒数作权重得到 $1,1/2,1/4$，归一化后为 $4/7,2/7,1/7$。回归时用这些权重乘以邻居目标再相加；分类时按类别累加权重后比较。距离为零时不能直接除零，程序需作特殊处理。\texttt{weights="uniform"}与\texttt{weights="distance"}对应这两类规则。

邻居数较小时，预测更多地由局部记录决定，容易受到个别噪声样本影响；邻居数较大时，预测更加平滑，却可能抹去小型结构。最合适的邻居数必须依据开发数据验证，而不能固定说“3一定比5好”。另外，均匀投票与距离加权会改变靠近类别交界处的行为。

\fig{ch07_knn_boundaries.pdf}{Iris 两个特征下不同邻居数得到的二维边界。图用于观察局部与平滑的变化，不是全特征的最终性能比较。}

\subsection{邻居数控制局部性与平滑程度}
$k$很小时，预测受少数记录支配，能跟随细节，也容易受测量噪声或错误标签影响；$k$较大时，平均或投票更稳定，却可能把跨越真实边界的记录混在一起。把$k$增加到全部训练样本，分类接近整体多数类，回归成为全训练目标均值，输入的局部差异被完全忽略。

因此$k$不是“越大越准确”的数量参数，而是模型平滑程度的控制量。固定$k$时，样本密集区域的邻域半径较小，稀疏区域较大；同一个$k$并不表示在所有位置使用同样的物理范围。特别在训练边界附近，最近邻仍可能很远，需要检查距离而不是只输出类别。

\subsection{分类边界与噪声的关系}
一近邻分类将特征空间分成若干区域，每个区域由离它最近的训练样本支配。新增或移动一条训练记录，会改变附近区域；错误标签可能形成一个小的错误“岛”。增加$k$通常减少这种细小结构，但也可能抹掉真正稀有的局部群体。

在二分类中取奇数$k$可以减少均匀投票的某类平票，但不能消除多分类平票，也不能消除恰好等距离的邻居边界。软件必须采用确定的平票处理，研究者不应把由编码顺序决定的结果解释成物理偏好。实际比较可记录平票或临界样本，尤其是小数据集。

\subsection{计算量也是方法选择的一部分}
直接查询需要对每个新样本与训练记录计算距离。若有$n$条训练、$q$条查询、$p$个特征，朴素距离计算量约随$qnp$增长，还要寻找最近的$k$个。树形索引可以在部分低维问题中加速，高维或特殊距离下未必有同样收益。

这说明KNN的训练时间短，不意味着部署预测也便宜。对于大规模天体目录，可以先在小样本上评估效果，再测量批量预测时间和内存。采用何种加速方法不属于本章核心，但应让学生知道计算成本发生在什么阶段。

\section{标准化与特征选择}
\texttt{KNeighborsClassifier} 和 \texttt{KNeighborsRegressor} 可直接使用 sklearn。为了公平比较，标准化需要与分类器一起放在 \texttt{Pipeline} 中，并在每个训练折内拟合。特征选择同样要遵守这一规则，否则观察到的验证表现可能包含提前使用验证信息的偏差。

\begin{lstlisting}[caption={Wine 上对照邻居数与标准化}]
from sklearn.datasets import load_wine
from sklearn.neighbors import KNeighborsClassifier
from sklearn.preprocessing import StandardScaler
from sklearn.pipeline import make_pipeline
from sklearn.model_selection import StratifiedKFold, cross_val_score

X, y = load_wine(return_X_y=True)
from sklearn.model_selection import train_test_split
X, X_hold, y, y_hold = train_test_split(
    X, y, test_size=0.2, stratify=y, random_state=42
)
# All comparisons below use development X, y only.
cv = StratifiedKFold(5, shuffle=True, random_state=7)
for k in [1, 3, 5, 11, 21]:
    raw = KNeighborsClassifier(n_neighbors=k)
    scaled = make_pipeline(StandardScaler(),
                           KNeighborsClassifier(n_neighbors=k))
    print(k, cross_val_score(raw, X, y, cv=cv).mean(),
          cross_val_score(scaled, X, y, cv=cv).mean())
\end{lstlisting}

\subsection{无关特征为什么会伤害距离}
假设两个特征足以区分类别，再加入100列与类别无关的随机测量。每一列都对距离平方贡献一个非负项，大量无关差异可能掩盖原来两列中的信号。把所有列标准化并不能自动识别哪些列无关，它只统一尺度。

一个简单计算可以说明维数效应。若两个独立标准高斯点的各维差独立、方差为2，则距离平方的均值为$2p$、方差为$8p$。相对标准差为$\sqrt{8p}/(2p)=\sqrt{2/p}$，随维数增加而减小。这是特定分布下距离趋于相对集中的例子，不是所有高维数据的普遍定量公式。

\fig{v3_ch07_concentration.pdf}{独立高斯人工数据中，两点距离的相对分散程度随维数变化。它展示一种高维几何现象，不是对全部科学数据的无条件结论。}

真实数据可能集中在低维结构上，所以不能仅凭列数大就断言KNN必然失败。应通过加入无关列的受控实验、特征选择和降维对照，检查问题来自维数、尺度还是样本密度。任何用标签进行的特征选择都要在训练折内完成。

\subsection{相关特征相当于重复强调某个方向}
若某一测量被复制成两列，欧氏距离中相应差平方被加了两次，其权重增加，而没有获得新的独立信息。两个高度相关但不完全相同的字段也可能产生类似效果。数据字典与相关性检查因此不仅服务于解释，还会实际改变基于距离的方法。

在专业任务中，有时这种权重是合理的，例如更重视某种已知物理特征；关键是明确选择，而不是因为文件恰好包含更多重复列就无意增加权重。

\section{鸢尾花与葡萄酒：尺度和邻居数对照}
课堂先在 Iris 只选两个花部特征，手算几个样本到新点的距离，找出最近的3条记录，再用 sklearn 校验。第二步使用全部四个 Iris 特征，检查二维直觉能否迁移。第三步换成 Wine，比较是否标准化以及不同 $k$ 选择带来的验证变化。

\fig{ch07_wine_k.pdf}{Wine 数据上邻居数和尺度处理的交叉验证对照。曲线来自固定折划分，仅用于说明本次实验中的参数影响。}

本章会发现一个常见现象：同一个算法在不同数据尺度和特征集合上可能呈现明显不同的表现。算法名称并不决定科学有效性，数据的表示方式也是模型的一部分。对不平衡类别或稀有对象，普通准确率不一定与检出需求一致，需要回到第4、6章的评价方法。

\subsection{实际五折结果及其解释}
随书实验在葡萄酒142条开发记录上使用完全相同的五折，比较原尺度与训练折内标准化。取$k=1$时，原尺度平均准确率约0.6756，标准化后约0.9433；取$k=9$时分别约0.6751和0.9648。本次候选中，标准化$k=25$的平均值约0.9788，但仍需保留其折间波动，而不是只报告最高小数。

\fig{v3_ch07_scale.pdf}{葡萄酒开发数据上，同一组折中的尺度和邻居数比较。误差棒是折间标准差，不是独立测试置信区间。}

这一结果支持“原始尺度在本数据的距离中不合适”，而不是“任何数据标准化都一定提高准确率”。特征本身的物理尺度有时已经具有合理权重。下一步可以固定标准化方式，只改变距离或邻居权重，每次比较限制变化因素。

\subsection{预测错误要回到邻居记录}
对一条误判样本，除了打印预测类别，还应查出邻居的原行号、距离、类别和关键特征。可能发现目标类本来与邻居重叠，或者查询点远离全部训练记录，也可能发现某一大尺度列主导距离。这样的诊断直接对应算法机制，优于只写“模型能力不足”。

\begin{lstlisting}[caption={在已经拟合的缩放器与KNN中追踪邻居}]
z_query = scaler.transform(query)
distances, positions = knn.kneighbors(z_query, n_neighbors=5)
neighbor_rows = X_train.iloc[positions[0]].copy()
neighbor_rows["distance"] = distances[0]
neighbor_rows["reference_class"] = y_train.iloc[positions[0]].to_numpy()
print(neighbor_rows)
\end{lstlisting}
这里\texttt{scaler}必须是在相同训练数据上拟合并用于训练该KNN的变换器。邻居位置是训练矩阵内的位置，不一定等于原始数据行号；因此保留索引有助于追踪。

\section{专业特征中的相似性}
若使用 Sonar 的60个频带特征，可先检查各频带是否具有相同含义和尺度，并比较 KNN 与逻辑回归。对于天体测光，星等、颜色及其不确定性可能具有不同的数值范围；如果观测数据横跨不同质量等级，最近邻关系也可能混合真实物理相似性与噪声相似性。

高维空间里，距离可能变得不够有区分度：当大量无关或噪声特征进入距离计算，真正有用的少量差异容易被淹没。第13章学习 PCA 后，可以再次检查低维表示上的近邻结果，但 PCA 在监督验证过程中必须放在训练折内部拟合。

\subsection{天体相似与观测相似可能混在一起}
两个天体的测光特征相近，可能因为物理性质相似，也可能因为共同消光、距离或仪器响应。用标准化距离寻找邻居时，模型只知道给定列的数值，没有自动区分这些来源。专业解释需要结合未进入距离的测量质量和背景信息。

例如先比较颜色而不是绝对亮度，可能减弱某些共同尺度影响，但同时丢掉亮度信息；是否合适取决于目标。对于物理信号，幅度归一化可能强调波形形状，却消除强度差异。每一种预处理都是关于“哪些差异应该保留”的判断，必须写进实验说明。

\subsection{什么时候不应强行给出预测}
KNN总能在非空训练集中找到最近记录，即使所有邻居都很远。因此“找到了邻居”并不意味着训练数据覆盖该样本。可以把最近邻距离、邻居类别一致性和测量质量作为诊断量；若超过由开发数据设定的适用范围，可以标记需要人工复核，而不是输出没有限定的高置信类别。

拒绝或复核规则本身也需要验证，不能事后只把预测错误的样本都标为不适用。应先说明阈值来源，并报告接受率与接受样本上的表现。这样才能看见提高可靠性所付出的覆盖范围代价。

\chapterpractice{使用 Iris 两个特征观察 $k=1,3,11$ 的决策边界，再在 Wine 全部特征下比较未标准化和标准化的交叉验证分数。最后在一份专业 CSV 中重新检查单位、变量范围以及样本独立性。}
\chapterexercise
\begin{enumerate}
\item 给出三条人工样本，手算欧氏距离与 $k=3$ 投票结果。
\item 若把某个长度变量的单位从米改为毫米，在未标准化的欧氏距离中会发生什么？
\item 为什么 KNN 的 \texttt{fit()} 和逻辑回归的 \texttt{fit()} 内部工作不同？
\item 若一个新对象在训练范围以外，最近的训练样本也可能很远。如何量化和报告这种情况？
\end{enumerate}

\section*{深化练习与方法检验}
本组题强调解释、计算与验证；D编号在本章内独立计数。计算题应写出所用定义和中间步骤，实验题保留实际输出，包括不支持预期的结果。
\begin{enumerate}[label=D\arabic*.]
\item 查询点为0，训练位置为0.1、0.2、2，类别依次1、0、0。比较3近邻均匀投票与按距离倒数加权的结果；再说明查询点恰等于训练点时为何需要特殊处理。
\item 两个邻居目标分别为2与8，距离为1与2。求均匀和距离倒数加权的回归预测。该输出必定位于哪些数值范围内？
\item 用 $p$ 维独立标准正态点说明平方距离为何随维数增长：计算两点平方距离的均值和方差，再说明相对波动随 $p$ 的变化。
\item 标准化后准确率提高，是否证明每个输入都应等权进入距离？提出一个只增加无关噪声特征的对照实验，并说明验证步骤。
\end{enumerate}
